\documentclass[10pt,a4paper]{amsart}
\usepackage[margin=19mm]{geometry}
\usepackage{amsmath,amssymb,mathtools}
\usepackage[scr=rsfs,cal=euler]{mathalfa}
\usepackage{xfrac}
\usepackage[unicode,hidelinks,bookmarks=false]{hyperref}
\newcommand{\ie}{i.e.\ }
\newcommand{\cf}{cf.\ }
\newcommand{\resp}{respectively\ }
\newcommand{\Subsection}{\subsection}
\providecommand{\nobreakdash}{\nobreak\hbox{-}}
\setlength{\emergencystretch}{2em}
\multlinegap0pt
\usepackage{bm}
\usepackage{amssymb}
\usepackage{algorithm}\usepackage{algpseudocode}
\usepackage{booktabs}
\newtheorem{lemma}{Lemma}
\title{An adaptive perfectly matched layer finite element method for acoustic-elastic interaction in periodic structures}
\author{Sijia Li, Lei Lin, Junliang Lv}
\date{Source: arXiv:2602.09055v1 (CC BY 4.0)}
\begin{document}
\maketitle

\noindent\emph{Source and licence.} An adaptive perfectly matched layer finite element method for acoustic-elastic interaction in periodic structures. Version arXiv:2602.09055v1 (CC BY 4.0), distributed by arXiv under the Creative Commons Attribution 4.0 International licence, http://creativecommons.org/licenses/by/4.0/, as recorded in the arXiv licence metadata for this exact version and shown on the abstract page https://arxiv.org/abs/2602.09055v1. Full licence notice: \texttt{licenses/arxiv-2602.09055-CC-BY-4.0.txt}.\par\smallskip

\noindent\textit{Editorial scope.} Counted unit: the article's lemma errorformula (source0.tex lines 873--881). The article's complete original proof of it is delivered with it (source0.tex lines 882--959, one \texttt{proof} environment of 78 lines). No mathematics is written, repaired or re-derived: the statement and the proof are the article's own lines, in the article's own notation, and no retained line is substituted or re-flowed.  Retained uncounted as the source-local context the proof needs: lines 387--389; lines 391--396; lines 510--512; lines 669--672; lines 674--679; lines 800--802; lines 809--816; lines 818--824.  Nothing was omitted from any retained span, so the package carries no omission record.  No retained line contains a citation, so the wrapper carries no bibliography and no bibliography entry.  No retained line contains a figure environment, an image-inclusion command or drawing code, so no image is delivered and the package declares no asset dependency.  Editorial formatting only, in addition: an AMS \texttt{amsart} preamble that restates the article's own declarations and macros used by the retained lines, namely the environments \texttt{lemma} on separate counters, together with the standard packages those lines need.  The retained environments print in this document as Lemma 1 in order of appearance, whereas the article prints the same items with the numbering of its own full document.  The complete native source of the article as distributed in the arXiv e-print for this exact version is bundled unchanged as \texttt{sources/194289/source0.tex}.
\begin{align}\label{DtN varproblem}
	A(\boldsymbol{U}, \boldsymbol{V})=L(\boldsymbol{V}),\quad \forall~ \boldsymbol{V}=(\varphi, \boldsymbol{\psi}) \in \mathscr{H}^1_\mathrm{qp}(\Omega).
\end{align}

\begin{align}\label{DtN varleft}
	A(\boldsymbol{U}, \boldsymbol{V})
	=\widetilde{A}(\boldsymbol{U}, \boldsymbol{V})
	+ \widetilde{B}_1(\boldsymbol{U}, \boldsymbol{V})
	+ \widetilde{B}_2(\boldsymbol{U}, \boldsymbol{V}),
\end{align}

\begin{align}\label{trun PMLvarproblem}
	B_\mathrm{D}(\hat{\boldsymbol{U}}, \boldsymbol{V})=L(\boldsymbol{V}), \quad \forall~ \boldsymbol{V}=(\varphi, \boldsymbol{\psi}) \in \mathscr{H}_{0,\mathrm{qp}}^1(\mathrm{D)}.
\end{align}

\begin{align}\label{PML varproblem}
	A^{\mathrm{PML}}(\boldsymbol{U}^{\mathrm{PML}}, \boldsymbol{V})
	=L(\boldsymbol{V}), \quad \forall ~ \boldsymbol{V}=(\varphi, \boldsymbol{\psi}) \in \mathscr{H}^1_\mathrm{qp}(\Omega),
\end{align}

\begin{align}\label{PML varleft}
	A^{\mathrm{PML}}(\boldsymbol{U}^{\mathrm{PML}}, \boldsymbol{V})
	=\widetilde{A}^{\mathrm{PML}}(\boldsymbol{U}^{\mathrm{PML}}, \boldsymbol{V})
	+\widetilde{B}_1^{\mathrm{PML}}(\boldsymbol{U}^{\mathrm{PML}}, \boldsymbol{V})
	+\widetilde{B}_2^{\mathrm{PML}}(\boldsymbol{U}^{\mathrm{PML}}, \boldsymbol{V})
\end{align}

\begin{align}\label{trun APPLvarproblem}
	B_\mathrm{D}(\hat{\boldsymbol{U}}_h, \boldsymbol{V}_h)=L(\boldsymbol{V}_h), \quad \forall~\boldsymbol{V}_h=(\varphi_h, \boldsymbol{\psi}_h) \in \stackrel{\circ}{\mathscr{H}_h}(\mathrm{D}).
\end{align}

\begin{align}\label{extenstion}
	\Delta_{\hat{\boldsymbol{x}}}  \overline{\tilde{\varphi}} + \kappa^2 \overline{\tilde{\varphi}}=0 & \quad\text { in } \Omega_+^{\mathrm{PML}}, \notag\\
	\Delta^*_{\hat{\boldsymbol{x}}}\overline{\tilde{\boldsymbol{\psi}}} +  \omega^2\rho \overline{\tilde{\boldsymbol{\psi}}} = 0 &\quad \text { in } \Omega^{\mathrm{PML}}_-, \notag\\
	\tilde{\varphi}\left(x_1, h_1\right)=\varphi\left(x_1, h_1\right) & \quad \text { on } \Gamma_+, \\
	\tilde{\boldsymbol{\psi}}\left(x_1, h_2\right)=\boldsymbol{\psi} \left(x_1, h_2\right) & \quad\text { on } \Gamma_-, \notag\\
	\tilde{\varphi}\left(x_1, h_1+\delta_1\right)=0 &\quad \text { on } \Gamma_+^{\mathrm{PML}},\notag \\
	\tilde{\boldsymbol{\psi}}\left(x_1, h_2-\delta_2\right)=0 &\quad \text { on } \Gamma_-^{\mathrm{PML}}.\notag
\end{align}

\begin{lemma}\label{lemma5}
	For any $(p,\boldsymbol{u}) \in \mathscr{H}^1_\mathrm{qp}\left(\Omega\right)$, $(\varphi, \boldsymbol{\psi}) \in \mathscr{H}^1_\mathrm{qp}\left(\Omega\right)$, one can obtain
	\begin{align*}
		\int_{\Gamma_+} \mathscr{T}_+^{\mathrm{PML}} p  \overline{\varphi} \mathrm{~d} s &=\int_{\Gamma_+} p  \partial_{x_2} \overline{\tilde{\varphi}} \mathrm{~d} s,\\
		\int_{\Gamma_-} \mathscr{T}_-^{\mathrm{PML}} \boldsymbol{u} \cdot \overline{\boldsymbol{\psi}} \mathrm{~d} s &=\int_{\Gamma_-} \boldsymbol{u} \cdot \mathscr{B} \overline{\tilde{\boldsymbol{\psi}}} \mathrm{~d} s.
	\end{align*}
\end{lemma}

\begin{lemma}\label{errorformula}
	For any $\boldsymbol{V}=(\varphi,\boldsymbol{\psi}) \in \mathscr{H}^1_\mathrm{qp}\left(\Omega\right)$, which can be extended to be a function in $\mathscr{H}_{0,\mathrm{qp}}^1\left(\mathrm{D}\right)$. Based on (\ref{extenstion}), and $\boldsymbol{V}_h=(\varphi_h,\boldsymbol{\psi}_h) \in \stackrel{\circ}{\mathscr{H}_h}(\mathrm{D})$, there holds
	\begin{align}\label{lemm6}
		A\left(\boldsymbol{U}-\hat{\boldsymbol{U}}_{h}, \boldsymbol{V}\right)=&~ L(\boldsymbol{V}-\boldsymbol{V}_h)
		-B_\mathrm{D}\left(\hat{\boldsymbol{U}}_h,\boldsymbol{V}-\boldsymbol{V}_h\right) \notag \\
		&+ \int_{\Gamma_+}\left(\mathscr{T}_+-\mathscr{T}_+^{\mathrm{PML}}\right)\hat{p}^{\mathrm{sc}}_h \overline{\varphi} \mathrm{~d} s
		+\int_{\Gamma_-}\left(\mathscr{T}_--\mathscr{T}_-^{\mathrm{PML}}\right)\hat{\boldsymbol{u}}_h \cdot \overline{\boldsymbol{\psi}} \mathrm{~d} s.
	\end{align}
\end{lemma}

\begin{proof}
	By (\ref{DtN varproblem}), (\ref{DtN varleft}), (\ref{PML varproblem}) and (\ref{PML varleft}), one can get
	\begin{align}
		A\left(\boldsymbol{U}-\hat{\boldsymbol{U}}_{h}, \boldsymbol{V}\right)=&~A\left(\boldsymbol{U}-\hat{\boldsymbol{U}}, \boldsymbol{V}\right)
		+ A\left(\hat{\boldsymbol{U}}-\hat{\boldsymbol{U}}_h, \boldsymbol{V}\right) \notag\\
		=&~A\left(\boldsymbol{U}-\hat{\boldsymbol{U}}, \boldsymbol{V}\right)
		+ A\left(\hat{\boldsymbol{U}}-\hat{\boldsymbol{U}}_h, \boldsymbol{V}\right) 
		-A^{\mathrm{PML}}\left(\hat{\boldsymbol{U}}-\hat{\boldsymbol{U}}_h, \boldsymbol{V}\right)
		+A^{\mathrm{PML}}\left(\hat{\boldsymbol{U}}-\hat{\boldsymbol{U}}_h, \boldsymbol{V}\right) \notag\\
		=&~A\left(\boldsymbol{U},\boldsymbol{V}\right)-A\left(\hat{\boldsymbol{U}},\boldsymbol{V}\right)+A^{\mathrm{PML}}\left(\hat{\boldsymbol{U}}-\hat{\boldsymbol{U}}_h, \boldsymbol{V}\right) \notag\\
		&-\int_{\Gamma_+}\left(\mathscr{T}_+-\mathscr{T}_+^{\mathrm{PML}}\right)\left(\hat{p}^{\mathrm{sc}}-\hat{p}^{\mathrm{sc}}_h\right) \overline{\varphi} \mathrm{~d}s -\int_{\Gamma_-}\left(\mathscr{T}_--\mathscr{T}_-^{\mathrm{PML}}\right)\left(\hat{\boldsymbol{u}}-\hat{\boldsymbol{u}}_h\right) \cdot \overline{\boldsymbol{\psi}} \mathrm{~d} s\notag \\
		=&~A\left(\boldsymbol{U},\boldsymbol{V}\right)-A\left(\hat{\boldsymbol{U}},\boldsymbol{V}\right)+A^{\mathrm{PML}}\left(\hat{\boldsymbol{U}},\boldsymbol{V}\right)-A^{\mathrm{PML}}\left(\hat{\boldsymbol{U}}, \boldsymbol{V}\right)
		+A^{\mathrm{PML}}\left(\hat{\boldsymbol{U}}-\hat{\boldsymbol{U}}_h, \boldsymbol{V}\right) \notag\\
		&-\int_{\Gamma_+}\left(\mathscr{T}_+-\mathscr{T}_+^{\mathrm{PML}}\right)\left(\hat{p}^{\mathrm{sc}}-\hat{p}^{\mathrm{sc}}_h\right) \overline{\varphi} \mathrm{~d} s
		-\int_{\Gamma_-}\left(\mathscr{T}_--\mathscr{T}_-^{\mathrm{PML}}\right)\left(\hat{\boldsymbol{u}}-\hat{\boldsymbol{u}}_h\right) \cdot \overline{\boldsymbol{\psi}} \mathrm{~d} s \notag\\
		=&\int_{\Gamma_+}\left(\mathscr{T}_+-\mathscr{T}_+^{\mathrm{PML}}\right)\hat{p}^{\mathrm{sc}} \overline{\varphi} \mathrm{~d} s
		+\int_{\Gamma_-}\left(\mathscr{T}_--\mathscr{T}_-^{\mathrm{PML}}\right)\hat{\boldsymbol{u}} \cdot \overline{\boldsymbol{\psi}} \mathrm{~d} s +A^{\mathrm{PML}}\left(\hat{\boldsymbol{U}}-\hat{\boldsymbol{U}}_h, \boldsymbol{V}\right) \notag\\
		&-\int_{\Gamma_+}\left(\mathscr{T}_+-\mathscr{T}_+^{\mathrm{PML}}\right)\left(\hat{p}^{\mathrm{sc}}-\hat{p}^{\mathrm{sc}}_h\right) \overline{\varphi} \mathrm{~d} s
		-\int_{\Gamma_-}\left(\mathscr{T}_--\mathscr{T}_-^{\mathrm{PML}}\right)\left(\hat{\boldsymbol{u}}-\hat{\boldsymbol{u}}_h\right) \cdot \overline{\boldsymbol{\psi}} \mathrm{~d} s \notag\\
		=&\int_{\Gamma_+}\left(\mathscr{T}_+-\mathscr{T}_+^{\mathrm{PML}}\right)\hat{p}^{\mathrm{sc}}_h \overline{\varphi} \mathrm{~d} s
		+\int_{\Gamma_-}\left(\mathscr{T}_--\mathscr{T}_-^{\mathrm{PML}}\right)\hat{\boldsymbol{u}}_h \cdot \overline{\boldsymbol{\psi}} \mathrm{~d} s +A^{\mathrm{PML}}\left(\hat{\boldsymbol{U}}-\hat{\boldsymbol{U}}_h, \boldsymbol{V}\right).\label{4.2.1}
	\end{align}
	Combining equation (\ref{PML varleft}) and Lemma \ref{lemma5} yields
	\begin{align*}
		A^{\mathrm{PML}}\left(\hat{\boldsymbol{U}}-\hat{\boldsymbol{U}}_h,\boldsymbol{V}\right)
		=&~B_{\Omega}\left(\hat{\boldsymbol{U}}-\hat{\boldsymbol{U}}_h, \boldsymbol{V}\right) -\int_{\Gamma_+}\mathscr{T}_+^{\mathrm{PML}}\left(\hat{p}^{\mathrm{sc}}-\hat{p}^{\mathrm{sc}}_h\right) \overline{\varphi} \mathrm{~d} s
		-\int_{\Gamma_-}\mathscr{T}_-^{\mathrm{PML}}\left(\hat{\boldsymbol{u}}-\hat{\boldsymbol{u}}_h\right) \cdot \overline{\boldsymbol{\psi}} \mathrm{~d} s\\
		=&~B_{\Omega}\left(\hat{\boldsymbol{U}}-\hat{\boldsymbol{U}}_h, \boldsymbol{V}\right) 
		- \int_{\Gamma_+}\left(\hat{p}^{\mathrm{sc}}-\hat{p}^{\mathrm{sc}}_h\right) \partial_{x_2} \overline{\varphi} \mathrm{~d} s
		- \int_{\Gamma_-} \left(\hat{\boldsymbol{u}}-\hat{\boldsymbol{u}}_h\right) \cdot \mathscr{B} \overline{\boldsymbol{\psi}} \mathrm{~d} s.
	\end{align*}
	Due to $\mathcal{L}_1 \overline{\varphi}=0$ in $\Omega^{\mathrm{PML}}_+$ and $\mathcal{L}_2 \overline{\boldsymbol{\psi}}=0$ in $\Omega^{\mathrm{PML}}_-$, it follows from the Green's formula and the first Betti's formula that
	\begin{align*}
		0 =& -\int_{\Omega^{\mathrm{PML}}_+}	\mathcal{L}_1\overline{\varphi}\left(\hat{p}^{\mathrm{sc}}-\hat{p}^{\mathrm{sc}}_h\right) \mathrm{~d} \boldsymbol{x} \notag\\
		=& \int_{\Omega^{\mathrm{PML}}_+}\left[\mathbb{A}\nabla \overline{\varphi} \cdot \nabla  \left(\hat{p}^{\mathrm{sc}}-\hat{p}^{\mathrm{sc}}_h\right)   
		- \kappa^2 s(x_2) \overline{\varphi}\left(\hat{p}^{\mathrm{sc}}-\hat{p}^{\mathrm{sc}}_h\right)\right] \mathrm{~d} \boldsymbol{x} 
		- \int_{\Gamma_+^\mathrm{PML}}\left(\hat{p}^{\mathrm{sc}}-\hat{p}^{\mathrm{sc}}_h\right) \partial_{\hat{x}_2} \overline{\varphi} \mathrm{~d} s
		+ \int_{\Gamma_+}\left(\hat{p}^{\mathrm{sc}}-\hat{p}^{\mathrm{sc}}_h\right) \partial_{x_2} \overline{\varphi} \mathrm{~d} s
	\end{align*}
	and 
	\begin{align*}
		0=&-\int_{\Omega^{\mathrm{PML}}_-}\mathcal{L}_2\overline{\boldsymbol{\psi}}\cdot\left(\hat{\boldsymbol{u}}-\hat{\boldsymbol{u}}_h\right)   \mathrm{~d} \boldsymbol{x} \notag\\
		= &\int_{\Omega^{\mathrm{PML}}_-}\left[\mathcal{S}_{\lambda, \mu}(\hat{\boldsymbol{u}}-\hat{\boldsymbol{u}}_h, \boldsymbol{\psi}) - \omega^2 s(x_2) \rho \left(\hat{\boldsymbol{u}}-\hat{\boldsymbol{u}}_h\right) \cdot \overline{\boldsymbol{\psi}}\right] \mathrm{~d} \boldsymbol{x}
		- \int_{\Gamma_-^\mathrm{PML}} \left(\hat{\boldsymbol{u}}-\hat{\boldsymbol{u}}_h\right) \cdot \mathscr{B}_{\hat{\boldsymbol{x}}} \overline{\boldsymbol{\psi}} \mathrm{~d} s
		+ \int_{\Gamma_-} \left(\hat{\boldsymbol{u}}-\hat{\boldsymbol{u}}_h\right) \cdot \mathscr{B} \overline{\boldsymbol{\psi}} \mathrm{~d} s.
	\end{align*}
	Furthermore, it follows from (\ref{trun PMLvarproblem}), (\ref{trun APPLvarproblem}) and the above two equalities that
	\begin{align}\label{4.2.2}
		A^{\mathrm{PML}}\left(\hat{\boldsymbol{U}}-\hat{\boldsymbol{U}}_h,\boldsymbol{V}\right)
		=&~ B_{\Omega}\left(\hat{\boldsymbol{U}}-\hat{\boldsymbol{U}}_h, \boldsymbol{V}\right) - \int_{\Gamma_+}\left(\hat{p}^{\mathrm{sc}}-\hat{p}^{\mathrm{sc}}_h\right) \partial_{x_2} \overline{\varphi} \mathrm{~d} s
		- \int_{\Gamma_-} \left(\hat{\boldsymbol{u}}-\hat{\boldsymbol{u}}_h\right) \cdot \mathscr{B} \overline{\boldsymbol{\psi}} \mathrm{~d} s \notag\\
		=&~B_\mathrm{D}\left(\hat{\boldsymbol{U}}-\hat{\boldsymbol{U}}_h, \boldsymbol{V}\right)
		- \int_{\Gamma_+^\mathrm{PML}}\left(\hat{p}^{\mathrm{sc}}-\hat{p}^{\mathrm{sc}}_h\right) \partial_{\hat{x}_2} \overline{\varphi} \mathrm{~d} s
		- \int_{\Gamma_-^\mathrm{PML}} \left(\hat{\boldsymbol{u}}-\hat{\boldsymbol{u}}_h\right) \cdot \mathscr{B}_{\hat{\boldsymbol{x}}} \overline{\boldsymbol{\psi}} \mathrm{~d} s\notag\\
		=&~ B_\mathrm{D}\left(\hat{\boldsymbol{U}}, \boldsymbol{V}\right) -B_\mathrm{D}\left(\hat{\boldsymbol{U}}_h, \boldsymbol{V}\right)
		+ B_\mathrm{D}\left(\hat{\boldsymbol{U}}_h, \boldsymbol{V}_h\right)
		- B_\mathrm{D}\left(\hat{\boldsymbol{U}}_h, \boldsymbol{V}_h\right)
		\notag\\
		&-\int_{\Gamma_+^\mathrm{PML}}\left(\hat{p}^{\mathrm{sc}}-\hat{p}^{\mathrm{sc}}_h\right) \partial_{\hat{x}_2} \overline{\varphi} \mathrm{~d} s
		- \int_{\Gamma_-^\mathrm{PML}} \left(\hat{\boldsymbol{u}}-\hat{\boldsymbol{u}}_h\right) \cdot \mathscr{B}_{\hat{\boldsymbol{x}}} \overline{\boldsymbol{\psi}} \mathrm{~d} s \notag\\
		=& ~L(\boldsymbol{V}) - B_\mathrm{D}\left(\hat{\boldsymbol{U}}_h, \boldsymbol{V}-\boldsymbol{V}_h\right)- L(\boldsymbol{V}_h)
		\notag\\
		&-\int_{\Gamma_+^\mathrm{PML}}\left(\hat{p}^{\mathrm{sc}}-\hat{p}^{\mathrm{sc}}_h\right) \partial_{\hat{x}_2} \overline{\varphi} \mathrm{~d} s
		- \int_{\Gamma_-^\mathrm{PML}} \left(\hat{\boldsymbol{u}}-\hat{\boldsymbol{u}}_h\right) \cdot \mathscr{B}_{\hat{\boldsymbol{x}}} \overline{\boldsymbol{\psi}} \mathrm{~d} s \notag\\
		=&~ L(\boldsymbol{V}-\boldsymbol{V}_h)
		-B_\mathrm{D}\left(\hat{\boldsymbol{U}}_h,\boldsymbol{V}-\boldsymbol{V}_h\right)
		-\int_{\Gamma_+^\mathrm{PML}}\left(\hat{p}^{\mathrm{sc}}-\hat{p}^{\mathrm{sc}}_h\right) \partial_{\hat{x}_2} \overline{\varphi} \mathrm{~d} s
		- \int_{\Gamma_-^\mathrm{PML}} \left(\hat{\boldsymbol{u}}-\hat{\boldsymbol{u}}_h\right) \cdot \mathscr{B}_{\hat{\boldsymbol{x}}} \overline{\boldsymbol{\psi}} \mathrm{~d} s.
	\end{align}
	Substituting (\ref{4.2.2}) into (\ref{4.2.1}) gives
	\begin{align*}
		A\left(\boldsymbol{U}-\hat{\boldsymbol{U}}_{h}, \boldsymbol{V}\right)=&~ L(\boldsymbol{V}-\boldsymbol{V}_h)
		-B_\mathrm{D}\left(\hat{\boldsymbol{U}}_h,\boldsymbol{V}-\boldsymbol{V}_h\right) \\
		&+ \int_{\Gamma_+}\left(\mathscr{T}_+-\mathscr{T}_+^{\mathrm{PML}}\right)\hat{p}^{\mathrm{sc}}_h \overline{\varphi} \mathrm{~d} s
		+\int_{\Gamma_-}\left(\mathscr{T}_s-\mathscr{T}_s^{\mathrm{PML}}\right)\hat{\boldsymbol{u}}_h \cdot \overline{\boldsymbol{\psi}} \mathrm{~d} s,
	\end{align*}
	where we have used $\hat{p}^{\mathrm{sc}}=\hat{p}^{\mathrm{sc}}_h=0$ on $\Gamma_+^{\mathrm{PML}}$ and $\hat{\boldsymbol{u}}=\hat{\boldsymbol{u}}_h=0$ on $\Gamma_-^{\mathrm{PML}}$. Thus the proof is completed.
\end{proof}

\end{document}
